% concatenated from scarduaSiegelgeometric_Pisa_submission.tex
\documentclass[11pt,reqno]{article}

\usepackage[utf8]{inputenc}
\usepackage[T1]{fontenc}
\usepackage{lmodern}
\usepackage{amsmath,amsthm,amsfonts,amssymb}
\usepackage{geometry}
\usepackage{enumitem}
\usepackage{hyperref}
\geometry{margin=1.08in}
\hypersetup{colorlinks=true,linkcolor=blue,citecolor=blue,urlcolor=blue}

\newtheorem{theorem}{Theorem}
\newtheorem{lemma}[theorem]{Lemma}
\newtheorem{corollary}{Corollary}
\newtheorem{proposition}[theorem]{Proposition}
\theoremstyle{remark}
\newtheorem*{remark}{Remark}
\newtheorem{exe}[theorem]{Example}
\newcommand{\C}{\mathbb C}
\newcommand{\R}{\mathbb R}
\newcommand{\D}{\mathbb D}
\newcommand{\F}{\mathcal F}
\newcommand{\G}{\mathcal G}
\newcommand{\Om}{\Omega}
\newcommand{\T}{\mathbb T}
\newcommand{\pd}{\partial}
\newcommand{\ol}{\overline}
\newcommand{\ip}[2]{\left\langle #1,#2\right\rangle}
\newcommand{\Sing}{\operatorname{Sing}}
\title{Geometric Orbital Linearization of Siegel Singularities in the Plane}
\author{Toshikazu Ito \and Bruno Sc\'ardua}
\date{}

\begin{document}
\maketitle

\begin{abstract}
We give geometric characterizations of analytically linearizable nondegenerate Siegel foliations in $\C^2$ using complex tangencies with shrinking strictly pseudoconvex hypersurfaces. Our principal result treats a fixed bounded smoothly bounded strictly pseudoconvex Reinhardt domain: in coordinates compatible with the eigendirections of the linear part, two-dimensional tangency loci along a shrinking sequence characterize analytic Siegel linearizability. For round spheres we prove a stronger statement: no a priori alignment of the Euclidean coordinate axes with the eigendirections is required; the tangency hypothesis itself forces the linear part to be unitarily diagonalizable with negative real eigenvalue ratio. We also establish a polynomial Shilov-boundary criterion on a round sphere and a logarithmic one-boundary criterion valid for arbitrary smooth strictly pseudoconvex domains, provided the embedded closed bidisc lies on the pseudoconvex side. Finally, we study the transverse-holomorphic dynamics of smooth tangency tori, including a rigidity theorem in the periodic case.
\end{abstract}
\medskip
\noindent\textbf{2020 Mathematics Subject Classification.} Primary 37F75; Secondary 32S65, 37F50, 32V40.
\tableofcontents
\section{Introduction and main results}

This work is about one of the main tools in theory of foliations and singularities, 
the effect of transverse sections in the dynamics and classification of the foliation.
In the real smooth framework this is quite classic and goes back to the works of Haefliger
(\cite{Haefliger}) and others. This line eventually led to the discovery of Reeb components and
finally to the celebrated theorem of Novikov about compact leaves in the 3-sphere (\cite{Camacho-LN}, “Novikov’s Theorem,” pp. 131–158). This was our original motivation for studying real transverse sections of holomorphic foliations. In the case of holomorphic foliations a natural question is about the existence of real hypersurfaces transverse to the foliation.  Let us be more precise.
For a smooth real hypersurface $H\subset\C^2$, a one-dimensional holomorphic foliation $\F$, and a regular point $p\in H$, we say that $\F$ is \emph{complex tangent} to $H$ at $p$ if the tangent line of $\F$ at $p$ is contained in the complex tangent line $T_p^{\C}H$. We denote the corresponding tangency set by $M(\F,H)$. A very first question is, for a holomorphic foliation in dimension two,  about the 
connection between the existence of   3-spheres transverse to the foliation and the description of the foliation inside the corresponding 4-ball. Denote by $S^3(R)\subset \mathbb C^2$ the radius $R>0$ round sphere centered at the origin. 

According to Ito-Douady, if $M(\F,S^3(R))=\emptyset$ for some $R>0$ then $\Sing(\F)\cap B^4(R)$ consists of a single point. After moving this singularity to the origin we can conclude that 
$M(\F,S^3(r))=\emptyset,\, \forall 0<r\leq R$ and the germ of $\F$ at the origin is nondegenerate and  in the Poincaré domain (\cite{Ito,Ito1996}).  Once we have this, Poincaré-Dulac theorem gives the possible analytic normal forms for the germ of $\F$ at the origin (\cite{Ilyashenko-Yakovenko}). 
Still owing to \cite{Ito}, the dynamics of $\F$ inside the 
ball $B(0;R)$ is well-known. In this work we investigate further the relation between the spherical tangency set and the analytic classification of the foliation.  If we consider a nondegenerate singularity, the remaining case is the Siegel case $\F: xdy - \lambda ydx + (...)=0, \lambda \in \mathbb R_-$. For the linear Siegel model the tangency set is a Clifford torus $T(r)=M(\F, S^3(R))$ (\S~\ref{section:siegellinear}). This is, in this rigid cartesian setting, a characterization of the linear  Siegel case (cf. Proposition~\ref{prop:cartesian}). 
For the general case of round spheres we have the following:

Given any fixed local complex coordinates $(z,w)$ centered at the origin, write
\[
 S^3_{(z,w)}(r)=\{(z,w)\in\C^2:|z|^2+|w|^2=r^2\}.
\]

\begin{theorem}[spherical geometric linearization]\label{thm:spherical-characterization}
Let $\F$ be a nondegenerate germ of one-dimensional holomorphic foliation at $0\in\C^2$. Then $\F$ is analytically linearizable of Siegel type if and only if there exist local holomorphic coordinates $(z,w)$ centered at the origin and a sequence $r_j\downarrow0$ such that
\[
 M\bigl(\F,S^3_{(z,w)}(r_j)\bigr)
\]
is a nonempty smooth real surface of dimension two for every $j$. Moreover, under the hypotheses above, in these Euclidean coordinates, after a unitary change of coordinates one has
\[
 \F:\quad \mu w\,dz+\lambda z\,dw=0,
 \qquad \lambda,\mu>0,
\]
and for every sufficiently small $r>0$ the tangency locus is the Clifford torus
\[
 |z|^2=\frac{\mu}{\lambda+\mu}r^2,
 \qquad
 |w|^2=\frac{\lambda}{\lambda+\mu}r^2.
\]
\end{theorem}



Since we are interested in classification under analytic changes of coordinates, we shall extend our framework by allowing tangency sets which are biholomorphic 
to the torus, not necessarily in the cartesian coordinates. 
More precisely, we study the analytic classification problem  where the tangency set is not-empty, biholomorphic to a torus in a natural sense.  
An important class of circled domains we consider  is the class of Reinhardt domains. 
A domain $D\subset\C^2$ is called \emph{Reinhardt} if it is invariant under the standard $\T^2$-action
\[
 (z,w)\longmapsto (e^{i\theta}z,e^{i\varphi}w),
 \qquad \theta,\varphi\in\R.
\]
Reinhardt symmetry alone does not imply pseudoconvexity. In our next  result we fix a bounded Reinhardt domain $D\Subset\C^2$ containing the origin and assume that its boundary is smooth and strictly pseudoconvex. Given $r>0$ we denote by $rD\subset \mathbb C^2$ its image under the homothety $(z,w)\longmapsto (rz,rw)$. 
\begin{theorem}[Reinhardt-domain geometric linearization]\label{thm:main}
Fix a bounded Reinhardt domain $D\Subset\C^2$ containing $0$, with smooth strictly pseudoconvex boundary. Let $\F$ be a nondegenerate germ of one-dimensional holomorphic foliation at $0\in\C^2$. The following conditions are equivalent:
\begin{enumerate}
\item There exists a system of local holomorphic coordinates $(z,w)$ centered at the origin and adapted to the eigendirections of the linear part of $\F$ (equivalently, the linear part is diagonal in these coordinates), together with a sequence $r_j\downarrow0$, such that, after identifying the coordinate neighborhood with a neighborhood of $0$ in $\C^2$, each boundary $\partial(r_jD)$ is contained in that neighborhood and
\[
 M\bigl(\F,\partial(r_jD)\bigr)
\]
is a nonempty smooth real surface of dimension two.
\item $\F$ is analytically linearizable of Siegel type; that is, there exists a local biholomorphism
\[
 \Psi:(\C^2,0)\longrightarrow(\C^2,0)
\]
such that, in the coordinates $(z,w)=\Psi(x,y)$,
\[
 \Psi_*\F:\quad z\,dw-\tau w\,dz=0,
 \qquad \tau\in\R_{<0}.
\]

\end{enumerate}
Moreover, whenever condition \textup{(1)} holds, in the same coordinates $(z,w)$ the foliation is already defined by
\[
 \mu w\,dz+\lambda z\,dw=0
\]
for some $\lambda,\mu>0$. For every sufficiently small $r>0$, the tangency locus $M(\F,\partial(rD))$ is the homothetic image of the unique Reinhardt torus on $\partial D$ tangent to the corresponding linear Siegel foliation.
\end{theorem}

The round sphere has an additional unitary symmetry, and this removes the spectral-alignment hypothesis.
Thus Theorem~\ref{thm:spherical-characterization} is not merely the specialization $D=\mathbb B^2$ of Theorem~\ref{thm:main}: its converse is stronger because the Euclidean axes need not initially be aligned with the eigendirections. The extra conclusion comes from the $U(2)$-symmetry of the sphere and from a simple rank-one tangency mechanism which will also reappear in the final section on smooth tangency tori.

\begin{remark}[strict pseudoconvexity without Reinhardt symmetry is not enough]
The Reinhardt hypothesis in Theorem~\ref{thm:main} is substantive. Consider the already linear Siegel foliation generated by
\[
 X_0=z\frac{\partial}{\partial z}-w\frac{\partial}{\partial w}
\]
and, for a small real $\varepsilon\neq0$, the domain
\[
 D_\varepsilon
 =\left\{(z,w):|z|^2+|w|^2+\varepsilon\Re z<1\right\}.
\]
This is a translated Euclidean ball, hence bounded with real-analytic strictly pseudoconvex boundary and containing $0$, but it is not Reinhardt. On $\partial D_\varepsilon$ the complex tangency equation is
\[
 |z|^2-|w|^2+\frac{\varepsilon}{2}z=0.
\]
Its imaginary part forces $\Im z=0$; writing $z=x\in\R$ and combining with the boundary equation gives
\[
 2x^2+\frac{3\varepsilon}{2}x=1.
\]
Thus the tangency locus is, when nonempty, a finite union of circles and is therefore one-dimensional. Since $X_0$ is homogeneous, the same conclusion holds on every homothetic boundary $\partial(rD_\varepsilon)$. Hence strict pseudoconvexity alone does not imply the forward direction of Theorem~\ref{thm:main}.
\end{remark}

In the resonant subcase, either of the shrinking-boundary theorems immediately yields a nonconstant holomorphic first integral; see Corollary~\ref{cor:resonant-first-integral}.

We also consider tangency sets arising as images of the distinguished boundary $\T^2$ of an analytically embedded bidisc. In the {\em Cartesian case} there is an elementary rigidity phenomenon: if $\Om=D_1\times D_2$ and $\partial D_1\times\partial D_2$ lies in a round sphere on which $\F$ is complex tangent along this product torus, then $D_1$ and $D_2$ are centered discs and the foliation is the corresponding linear Siegel foliation up to multiplication of a generator by a holomorphic factor. The precise statement and proof are given in Proposition~\ref{prop:cartesian}.

The polynomial one-boundary criterion below is genuinely spherical. 

\begin{theorem}[polynomial Shilov rigidity]\label{thm:polynomial}
Let 
\[
 \Phi: W \longrightarrow \Om\subset\C^2
\]
be a biholomorphic embedding induced by a polynomial map, defined in a neighborhood $W$ of $\ol{\D}^2$, with $\Phi(0)=0$, and set
\[
 T=\Phi(\T^2),\qquad \T^2=\{(z,w):|z|=|w|=1\}.
\]
Assume that $T\subset S^3(R)$, that $\F$ is holomorphic in a neighborhood of $\ol\Om$, that
\[
 T\subset M(\F,S^3(R)),
\]
and that $0=\Phi(0)$ is an isolated nondegenerate Siegel singularity of $\F$.
 Then $\Phi$ is linear, up to a unitary change of target coordinates, and the germ of $\F$ at $0$ is analytically linearizable as a Siegel type singularity.
\end{theorem}

The logarithmic criterion below is more flexible and does not require a sphere, or even Reinhardt symmetry. 

\begin{theorem}[logarithmic rigidity for strictly pseudoconvex boundaries]\label{thm:logarithmic}
Let $D\Subset\C^2$ be a bounded domain with smooth strictly pseudoconvex boundary, and let
\[
 \Phi:\D^2\longrightarrow \Om\subset\C^2
\]
be a biholomorphic embedding extending holomorphically to a neighborhood of $\ol{\D}^2$, with $\Phi(0)=0$, such that
\[
 \Phi(\ol{\D}^2)\subset\ol D,
 \qquad
 T:=\Phi(\T^2)\subset\partial D.
\]
Let $\F$ be a one-dimensional holomorphic foliation defined in a neighborhood of $\Phi(\ol{\D}^2)$ and assume
\[
 T\subset M(\F,\partial D).
\]
Suppose that the pulled-back foliation $\G=\Phi^*\F$ is defined on a neighborhood of $\ol{\D}^2$ by
\[
 \omega=a(z,w)\frac{dz}{z}+b(z,w)\frac{dw}{w},
\]
where $a$ and $b$ are holomorphic units. Then the germ of $\F$ at $0$ is analytically linearizable of Siegel type.
\end{theorem}

The four principal rigidity results have complementary geometric characters. Theorem~\ref{thm:main} treats shrinking strictly pseudoconvex Reinhardt boundaries under compatibility with the eigendirections. Theorem~\ref{thm:spherical-characterization} uses the larger symmetry of the round sphere to remove that compatibility hypothesis and in fact forces unitary spectral alignment. Theorem~\ref{thm:polynomial} is a one-boundary spherical result in which polynomiality rigidifies the embedding itself. Theorem~\ref{thm:logarithmic}, by contrast, applies to an arbitrary smooth strictly pseudoconvex boundary once the embedded closed bidisc lies on the pseudoconvex side and the pulled-back foliation admits the global logarithmic representation above.

A second theme concerns the local dynamics near a smooth tangency torus. In the periodic case we prove that the singular characteristic line field admits a nonsingular real-analytic desingularization and that both the ambient holomorphic holonomy and the Poincar\'e return map of the desingularized flow are trivial; see Theorem~\ref{thm:periodic-torus-rigidity}. The differential identity underlying that desingularization is already used in the proof of Theorem~\ref{thm:spherical-characterization} to prevent a sequence of two-dimensional tangency surfaces from collapsing onto a lower-dimensional critical tangency set of the linear part.




\section{The spherical theorem and automatic spectral alignment}\label{sec:spherical}

We now prove Theorem~\ref{thm:spherical-characterization}. We first isolate the rank-one mechanism that prevents two-dimensional tangency surfaces from collapsing, after rescaling, onto a lower-dimensional critical zero set.

\begin{lemma}[rank-one spherical tangency]\label{lem:spherical-rank}
Let $X=P_1\partial_{z_1}+P_2\partial_{z_2}$ be holomorphic and regular near a point of $S^3(r)$, put
\[
 \rho=|z_1|^2+|z_2|^2-r^2,
 \qquad g=X(\rho)=P_1\ol z_1+P_2\ol z_2.
\]
At every tangency point $g=0$ one has
\[
 \ol X(g)=|P_1|^2+|P_2|^2=\|X\|^2>0.
\]
Consequently $dg|_{TS^3(r)}\ne0$. If $g^{-1}(0)\subset S^3(r)$ is a smooth real surface of dimension two, then the real map $g:S^3(r)\to\C$ has rank exactly one along that surface.
\end{lemma}

\begin{proof}
Holomorphicity of $P_1,P_2$ gives
\[
 \ol X(g)=|P_1|^2+|P_2|^2.
\]
At $g=0$, both real and imaginary parts of $X$ are tangent to the sphere. Since the displayed derivative is nonzero, the real differential of $g$ restricted to the tangent space of the sphere cannot vanish. A two-dimensional zero set inside the three-dimensional sphere has real codimension one, so its defining map cannot have rank two there; hence the rank is exactly one.
\end{proof}

\begin{lemma}[linear spherical alignment]\label{lem:linear-spherical-alignment}
Let $X_1(z)=Az$ be an invertible complex linear vector field on $\C^2$. If $M(X_1,S^3(1))$ contains a smooth real surface of dimension two, then $A$ is unitarily diagonalizable and its two eigenvalues have negative real ratio.
\end{lemma}

\begin{proof}
By a unitary change of coordinates, which preserves $S^3(1)$, put $A$ in Schur form
\[
 A=\begin{pmatrix}\alpha&\gamma\\0&\beta\end{pmatrix},
 \qquad \alpha\beta\ne0.
\]
Thus
\[
 X_1=(\alpha z+\gamma w)\frac{\pd}{\pd z}
      +\beta w\frac{\pd}{\pd w}.
\]
On $S^3(1)$ the tangency equation is
\begin{equation}\label{eq:schur-tangency}
 \alpha|z|^2+\beta|w|^2+\gamma w\ol z=0.
\end{equation}
Write $s=|z|^2$, $|w|^2=1-s$, and, when $0<s<1$,
\[
 w\ol z=\sqrt{s(1-s)}e^{i\theta}.
\]
If $\gamma\ne0$, equation~\eqref{eq:schur-tangency} determines
\[
 e^{i\theta}=-\frac{\alpha s+\beta(1-s)}{\gamma\sqrt{s(1-s)}},
\]
and therefore $s$ must satisfy
\begin{equation}\label{eq:schur-modulus}
 |\alpha s+\beta(1-s)|^2=|\gamma|^2s(1-s).
\end{equation}
The two sides give a real polynomial equation of degree at most two. It is not an identity, since at $s=0$ the difference equals $|\beta|^2>0$. Hence only finitely many values of $s$ occur; for each such value the relative phase is fixed, while the common Hopf phase is free. The tangency locus is therefore a finite union of circles, contradicting the existence of a two-dimensional surface. Thus $\gamma=0$.

Equation~\eqref{eq:schur-tangency} becomes
\[
 \alpha s+\beta(1-s)=0.
\]
A two-dimensional tangency component cannot lie on a coordinate circle, so $0<s<1$, and
\[
 \frac{\alpha}{\beta}=-\frac{1-s}{s}\in\R_{<0}.
\]
Thus the Schur form is diagonal and the eigenvalue ratio is negative real.
\end{proof}

\begin{lemma}[no collapse of spherical tangency surfaces]\label{lem:no-collapse}
Let $X=X_1+O(2)$ be a holomorphic vector-field germ with invertible linear part. Put
\[
 X_r(\zeta)=r^{-1}X(r\zeta),\qquad
 \rho(\zeta)=|\zeta_1|^2+|\zeta_2|^2-1,
 \qquad g_r=X_r(\rho)|_{S^3(1)}.
\]
If $r_j\downarrow0$ and each $g_{r_j}^{-1}(0)$ is a nonempty smooth real surface of dimension two, then $g_0^{-1}(0)=M(X_1,S^3(1))$ contains a smooth real surface of dimension two.
\end{lemma}

\begin{proof}
Choose $p_j\in g_{r_j}^{-1}(0)$ and, after passing to a subsequence, $p_j\to p\in S^3(1)$. Since $X_r\to X_1$ in $C^1$ on the sphere,
\[
 g_0(p)=0.
\]
By Lemma~\ref{lem:spherical-rank}, $dg_0|_{TS^3}(p)\ne0$. On the other hand, its real rank cannot be two. Indeed, if it were two, the submersion theorem and $C^1$ convergence would imply that $g_r^{-1}(0)$ is locally one-dimensional near $p$ for all sufficiently small $r$, contrary to the presence of the two-dimensional zero surface through $p_j$.
Hence
\[
 \operatorname{rank}_{\R}dg_0|_{TS^3}(p)=1.
\]
After multiplying all $g_r$ by one fixed unimodular constant, we may assume that
\[
 d(\Im g_0)|_{TS^3}(p)\ne0.
\]
Choose a small coordinate neighborhood $U$ of $p$ in $S^3(1)$ so that the implicit-function theorem represents, for every sufficiently small $r$, the set
\[
 \Sigma_r=\{\Im g_r=0\}\cap U
\]
as one connected real-analytic graph over a fixed two-disc. These graphs depend continuously in $C^1$ on $r$ and converge to $\Sigma_0$.

For large $j$, $p_j\in U$ and the two-dimensional tangency surface $T_j=g_{r_j}^{-1}(0)$ meets $U$. Since
\[
 T_j\cap U\subset\Sigma_{r_j}
\]
and both are smooth real surfaces of dimension two near $p_j$, the inclusion is locally open. Thus $\Re g_{r_j}$ vanishes on a nonempty open subset of the connected real-analytic surface $\Sigma_{r_j}$. By the real-analytic identity principle,
\[
 \Re g_{r_j}\equiv0\qquad\hbox{on }\Sigma_{r_j}.
\]
Passing to the limit in the fixed graph coordinates gives
\[
 g_0\equiv0\qquad\hbox{on }\Sigma_0.
\]
Hence $g_0^{-1}(0)$ contains the smooth real surface $\Sigma_0$ through $p$.
\end{proof}

\begin{proof}[Proof of Theorem~\ref{thm:spherical-characterization}]
Assume first that $\F$ is analytically linearizable of Siegel type. Choose linearizing coordinates $(u,v)$, so that a generator is
\[
 X_0=\lambda u\frac{\pd}{\pd u}-\mu v\frac{\pd}{\pd v},
 \qquad \lambda,\mu>0.
\]
Taking the Euclidean spheres in these coordinates gives, for every small $r$,
\[
 X_0(|u|^2+|v|^2)=\lambda|u|^2-\mu|v|^2,
\]
and therefore the tangency locus is the Clifford torus
\[
 |u|^2=\frac{\mu}{\lambda+\mu}r^2,
 \qquad
 |v|^2=\frac{\lambda}{\lambda+\mu}r^2.
\]
Thus the tangency condition holds.

Conversely, suppose the tangency condition holds in fixed local Euclidean coordinates $(z,w)$, without assuming that their axes are eigendirections. Choose a holomorphic generator
\[
 X=X_1+X_2+\cdots
\]
with invertible linear part. Rescaling the spheres to $S^3(1)$ gives the family $X_r=r^{-1}X(r\cdot)$. Lemma~\ref{lem:no-collapse} shows that the linear tangency locus $M(X_1,S^3(1))$ contains a smooth real surface of dimension two. Lemma~\ref{lem:linear-spherical-alignment} therefore implies that $X_1$ is unitarily diagonalizable and has negative real eigenvalue ratio. Performing this unitary change does not alter any Euclidean sphere or the dimension of any tangency locus. After multiplication by a nonzero constant we may consequently assume
\[
 X_1=X_0=\lambda z\frac{\pd}{\pd z}-\mu w\frac{\pd}{\pd w},
 \qquad \lambda,\mu>0.
\]

It remains to eliminate the nonlinear terms. Suppose $X\wedge X_0\not\equiv0$. Multiplication by a holomorphic unit does not change the foliation or its tangency locus. Removing the homogeneous orbital multiples of $X_0$ that precede the first nontrivial wedge degree, we may write, for some $m\ge2$,
\[
 X=X_0+X_m+O(m+1),
 \qquad
 X_m=P_m\frac{\pd}{\pd z}+Q_m\frac{\pd}{\pd w},
\]
where $P_m,Q_m$ are homogeneous of degree $m$ and $X_m$ is not a homogeneous multiple of $X_0$.

On $S^3(1)$ set
\[
 G_r=X_r(|z|^2+|w|^2-1).
\]
Then
\[
 G_r=G_0+r^{m-1}F_m+O(r^m),
 \qquad
 G_0=\lambda|z|^2-\mu|w|^2,
\]
with
\[
 F_m=P_m\ol z+Q_m\ol w.
\]
The equation $G_0=0$ cuts out the Clifford torus
\[
 T_0=\left\{|z|^2=\frac{\mu}{\lambda+\mu},\quad
 |w|^2=\frac{\lambda}{\lambda+\mu}\right\}
\]
transversely. Hence for all sufficiently small $r$, the real equation $\Re G_r=0$ is a unique connected graph $\Sigma_r\simeq\T^2$ over $T_0$, and there are no other solutions. For $r=r_j$, the actual tangency set $\{G_r=0\}$ is a nonempty compact smooth surface of dimension two contained in $\Sigma_r$; it is therefore open and closed in $\Sigma_r$, hence equals $\Sigma_r$. Dividing the imaginary part by $r_j^{m-1}$ and passing to the limit gives
\[
 \Im F_m=0\qquad\hbox{on }T_0.
\]

Every Fourier frequency occurring in $P_m\ol z+Q_m\ol w$ on $T_0$ has component sum $m-1>0$. A real-valued function must contain the opposite conjugate of each nonzero frequency, whose component sum would be $-(m-1)$; no such frequency occurs. Therefore
\[
 F_m\equiv0\qquad\hbox{on }T_0.
\]
Using
\[
 \ol z=\frac{\mu}{\lambda+\mu}\frac1z,
 \qquad
 \ol w=\frac{\lambda}{\lambda+\mu}\frac1w
\]
on $T_0$, we obtain
\[
 \mu wP_m+\lambda zQ_m\equiv0.
\]
Thus $X_m$ is a homogeneous holomorphic multiple of $X_0$, contradicting its choice. Hence $X\wedge X_0\equiv0$, and therefore
\[
 X=h(z,w)X_0
\]
for a holomorphic unit $h$. The foliation is analytically linear in the unitary coordinates. The explicit Clifford-torus formula above then holds for every sufficiently small sphere.
\end{proof}

\begin{corollary}[resonant subcase]\label{cor:resonant-first-integral}
Under the hypotheses of either Theorem~\ref{thm:main} or Theorem~\ref{thm:spherical-characterization}, if the resulting Siegel eigenvalue ratio is rational, then the germ of $\F$ admits a nonconstant holomorphic first integral.
\end{corollary}

\begin{proof}
By the corresponding theorem, the foliation germ is analytically linearizable. After a constant rescaling and analytic change of coordinates, a generator is orbitally equivalent to
\[
 pz\frac{\pd}{\pd z}-qw\frac{\pd}{\pd w},
 \qquad p,q\in\mathbb N,\quad \gcd(p,q)=1.
\]
The monomial
\[
 F(z,w)=z^q w^p
\]
is a holomorphic first integral of the linear model, and its pullback by the analytic conjugacy is a holomorphic first integral of the original germ.
\end{proof}












\section{Proof of Theorem~\ref{thm:main}}

We first isolate the geometry of a strictly pseudoconvex Reinhardt boundary. Put
\[
 s=|z|^2,\qquad t=|w|^2.
\]
Since $D$ is pseudoconvex Reinhardt and contains the origin, it is complete Reinhardt. Consequently the portion of $\partial D$ in $\{zw\neq0\}$ can be written as
\begin{equation}\label{eq:reinhardt-profile}
 t=\varphi(s),\qquad 0<s<a,
\end{equation}
where $\varphi(s)>0$, $\varphi(0)>0$, and $\varphi(s)\to0$ as $s\to a^-$. We may use
\[
 \rho(z,w)=|w|^2-\varphi(|z|^2)
\]
as a defining function along this part of the boundary.

\begin{lemma}[Levi monotonicity of the Reinhardt profile]\label{lem:reinhardt-profile}
With the notation above, the function
\[
 h(s)=-\frac{s\varphi'(s)}{\varphi(s)},\qquad 0<s<a,
\]
is strictly increasing and maps $(0,a)$ onto $(0,+\infty)$.
\end{lemma}

\begin{proof}
A complex tangent vector to $\rho=0$ is proportional to
\[
 \xi=(\ol w,\varphi'(s)\ol z).
\]
Since
\[
 \rho_{z\ol z}=-\varphi'(s)-s\varphi''(s),
 \qquad
 \rho_{w\ol w}=1,
\]
the Levi form along $\xi$ is
\[
 \mathcal L_\rho(\xi)
 =\varphi(s)\bigl[-\varphi'(s)-s\varphi''(s)\bigr]
   +s(\varphi'(s))^2.
\]
Direct differentiation gives the exact identity
\begin{equation}\label{eq:levi-h}
 \mathcal L_\rho(\xi)=\varphi(s)^2h'(s).
\end{equation}
Strict pseudoconvexity therefore gives $h'(s)>0$.

At the $w$-axis, smoothness gives bounded $\varphi'$ while $\varphi(0)>0$, and hence
\[
 \lim_{s\to0^+}h(s)=0.
\]
At the other endpoint, $\varphi(s)\to0$ as $s\to a^-$. If $h$ were bounded above by a constant $C$ near $a$, then
\[
 -\frac{\varphi'(s)}{\varphi(s)}\leq\frac{C}{s}.
\]
Integration on a terminal interval would give a positive lower bound for $\varphi(s)$ as $s\to a^-$, a contradiction. Thus
\[
 \lim_{s\to a^-}h(s)=+\infty.
\]
Together with strict monotonicity this proves the lemma.
\end{proof}

\begin{lemma}[fixed-coordinate Reinhardt rigidity]\label{lem:reinhardt-rigidity}
Let $X$ be a holomorphic vector-field germ at $0\in\C^2$ with isolated nondegenerate singularity, and suppose that in the chosen coordinates its linear part is diagonal,
\[
 X_1=\alpha z\frac{\pd}{\pd z}+\beta w\frac{\pd}{\pd w},
 \qquad \alpha\beta\ne0.
\]
Assume that there exists a sequence $r_j\downarrow0$ such that, for every $j$, the complex tangency locus $M(X,\partial(r_jD))$ is a nonempty smooth real surface of dimension two. Then
\[
 \frac{\alpha}{\beta}\in\R_{<0},
\]
and, after multiplication of $X$ by a nonzero constant,
\[
 X=h_0(z,w)\left(
 \lambda z\frac{\pd}{\pd z}-\mu w\frac{\pd}{\pd w}
 \right),
 \qquad \lambda,\mu>0,
\]
for a holomorphic unit $h_0$.
\end{lemma}

\begin{proof}
Rescale $z=r\zeta$, $w=r\omega$. The vector field
\[
 \widetilde X_r=r^{-1}X(r\zeta,r\omega)
\]
converges in $C^1$ on compact sets to $X_1$. Choose a tangency point on each $\partial(r_jD)$ and rescale it to $p_j\in\partial D$. Passing to a subsequence, $p_j\to p_0\in\partial D$, and
\[
 X_1(\rho)(p_0)=0
\]
for any local Reinhardt defining function $\rho$.

The point $p_0$ cannot lie on a coordinate axis. For example, at a point $(0,w)$ of $\partial D$ with $w\ne0$, Reinhardt invariance gives $\rho_z=0$, while regularity of the boundary gives $\rho_w\ne0$; hence $X_1(\rho)=\beta w\rho_w\ne0$. The argument on the other axis is identical. Thus $p_0$ has both coordinates nonzero. Writing locally $\rho=\rho(s,t)$, the complete Reinhardt property together with strict pseudoconvexity (equivalently, with the monotonicity of the profile established in Lemma~\ref{lem:reinhardt-profile}) gives $\rho_s,\rho_t>0$ along the positive part of the boundary, after choosing the outward orientation of $\rho$, and therefore
\[
 \alpha s\rho_s+\beta t\rho_t=0
\]
at $p_0$. It follows that
\[
 \frac{\alpha}{\beta}=-\frac{t\rho_t}{s\rho_s}\in\R_{<0}.
\]
After multiplication by a nonzero constant we may write
\[
 X_1=X_0=\lambda z\frac{\pd}{\pd z}-\mu w\frac{\pd}{\pd w},
 \qquad \lambda,\mu>0.
\]

For the profile~\eqref{eq:reinhardt-profile},
\begin{equation}\label{eq:linear-reinhardt-tangency}
 X_0(\rho)
 =-\lambda s\varphi'(s)-\mu\varphi(s)
 =\varphi(s)\bigl(\lambda h(s)-\mu\bigr).
\end{equation}
By Lemma~\ref{lem:reinhardt-profile}, this vanishes on exactly one Reinhardt torus
\[
 T_D=\{|z|^2=s_0,\ |w|^2=t_0\},
 \qquad t_0=\varphi(s_0),
\]
where $h(s_0)=\mu/\lambda$. The zero is transverse, since
\begin{equation}\label{eq:transverse-reinhardt-zero}
 \frac{d}{ds}X_0(\rho)\Big|_{s=s_0}
 =\lambda\varphi(s_0)h'(s_0)>0.
\end{equation}

Suppose now, toward a contradiction, that $X\wedge X_0\not\equiv0$. Multiplication by a holomorphic unit does not change the complex tangency set. We may therefore remove, one degree at a time, every homogeneous term which is a holomorphic multiple of $X_0$ before the first degree at which the wedge with $X_0$ is nonzero. Since only finitely many degrees precede that first degree, this requires only finitely many unit multiplications. Thus, for some $m\ge2$,
\[
 X=X_0+X_m+O(m+1),
 \qquad
 X_m=P_m(z,w)\frac{\pd}{\pd z}+Q_m(z,w)\frac{\pd}{\pd w},
\]
where $P_m,Q_m$ are homogeneous of degree $m$ and $X_m$ is not a homogeneous holomorphic multiple of $X_0$.

On the fixed boundary $\partial D$, after rescaling, put
\[
 G_r=\widetilde X_r(\rho).
\]
Then
\begin{equation}\label{eq:Gr-reinhardt}
 G_r=G_0+r^{m-1}F_m+O(r^m),
\end{equation}
where
\[
 G_0=X_0(\rho)
\]
and
\[
 F_m=\rho_sP_m\ol z+\rho_tQ_m\ol w.
\]
The convergence is uniform with first derivatives on $\partial D$.

By~\eqref{eq:linear-reinhardt-tangency}--\eqref{eq:transverse-reinhardt-zero}, $G_0$ has exactly one zero torus $T_D$ and is transverse to zero there. Choose a small Reinhardt tubular neighborhood $U$ of $T_D$. On the compact set $\partial D\setminus U$, $|G_0|$ is bounded below by a positive constant, so for sufficiently small $r$ the equation $\Re G_r=0$ has no solutions there. Inside $U$, the implicit-function theorem and the uniform nonvanishing radial derivative show that $\Re G_r=0$ is a unique connected global graph
\[
 \Sigma_r\simeq\T^2
\]
over $T_D$.

For $r=r_j$, the tangency locus
\[
 M_r=\{G_r=0\}
\]
is, by hypothesis, a nonempty smooth compact real surface of dimension two contained in the connected two-manifold $\Sigma_r$. Its inclusion is locally open, while compactness makes it closed. Hence
\[
 M_r=\Sigma_r.
\]
In particular, $\Im G_r=0$ on the whole graph. Dividing the imaginary part of~\eqref{eq:Gr-reinhardt} by $r_j^{m-1}$ and letting $j\to\infty$ gives
\[
 \Im F_m=0
\]
on $T_D$.

On $T_D$ the quantities
\[
 c_1=\rho_s(s_0,t_0),\qquad c_2=\rho_t(s_0,t_0)
\]
are positive constants, so
\[
 F_m=c_1P_m\ol z+c_2Q_m\ol w.
\]
Every Fourier frequency occurring in $P_m\ol z$ or $Q_m\ol w$ has component sum $m-1>0$. A real-valued function on $\T^2$ must contain with each frequency its opposite conjugate frequency, whose component sum would be $-(m-1)$. No such frequencies occur. Therefore
\[
 F_m\equiv0\qquad\text{on }T_D.
\]
Since on $T_D$
\[
 \ol z=\frac{s_0}{z},\qquad \ol w=\frac{t_0}{w},
\]
we obtain
\[
 c_1s_0\,wP_m+c_2t_0\,zQ_m=0
\]
on $T_D$, hence identically as a holomorphic polynomial. The linear tangency relation on $T_D$ is
\[
 \lambda c_1s_0=\mu c_2t_0,
\]
and consequently
\[
 \mu wP_m+\lambda zQ_m\equiv0.
\]
Since $z$ and $w$ are relatively prime,
\[
 P_m=zR_{m-1},\qquad
 Q_m=-\frac{\mu}{\lambda}wR_{m-1}
\]
for a homogeneous polynomial $R_{m-1}$. Thus $X_m=(R_{m-1}/\lambda)X_0$, contradicting the choice of $m$. Hence
\[
 X\wedge X_0\equiv0.
\]
Writing $X=P\partial_z+Q\partial_w$, this identity is
\[
 \mu wP+\lambda zQ=0.
\]
Putting successively $z=0$ and $w=0$ gives $z\mid P$ and $w\mid Q$, and therefore
\[
 X=h_0(z,w)X_0
\]
with $h_0$ holomorphic. Comparison of linear parts gives $h_0(0)=1$, so $h_0$ is a unit.
\end{proof}

\begin{proof}[Proof of Theorem~\ref{thm:main}]
Assume first that condition \textup{(2)} holds and choose analytic linearizing coordinates $(z,w)$. In those coordinates the foliation is generated by
\[
 X_0=\lambda z\frac{\pd}{\pd z}-\mu w\frac{\pd}{\pd w},
 \qquad \lambda,\mu>0.
\]
These coordinates are automatically adapted to the eigendirections of the linear part. Lemma~\ref{lem:reinhardt-profile} and equation~\eqref{eq:linear-reinhardt-tangency} show that $X_0$ has exactly one smooth Reinhardt tangency torus on $\partial D$. By homogeneity, its tangency locus on $\partial(rD)$ is exactly the homothetic torus $rT_D$ for every sufficiently small $r$. Thus \textup{(1)} holds.

Conversely, assume \textup{(1)} and use the coordinates occurring there. Choose a holomorphic generator $X$ of $\F$. Its linear part is diagonal, so Lemma~\ref{lem:reinhardt-rigidity} gives
\[
 X=h_0(z,w)\left(
 \lambda z\frac{\pd}{\pd z}-\mu w\frac{\pd}{\pd w}
 \right)
\]
with $h_0$ a holomorphic unit. Thus the foliation itself is linear Siegel in those coordinates, proving \textup{(2)} and the final assertion.
\end{proof}

\section{Cartesian one-sphere rigidity}

We next isolate the elementary Cartesian extreme of the one-sphere problem.

\begin{proposition}[Cartesian one-sphere rigidity]\label{prop:cartesian}
Let $\Om=D_1\times D_2\subset\C^2$ be a Cartesian product of bounded connected planar domains containing $0$, and suppose
\[
 \partial D_1\times\partial D_2\subset S^3(R).
\]
Let $\F$ be holomorphic near $\ol\Om$ and complex tangent to $S^3(R)$ along $\partial D_1\times\partial D_2$. Then there exist $a,b>0$, $a^2+b^2=R^2$, such that
\[
 D_1=\{|x|<a\},\qquad D_2=\{|y|<b\},
\]
and, for a local generator $X$ of $\F$ near $0$,
\[
 X=h(x,y)\left(b^2x\frac{\pd}{\pd x}-a^2y\frac{\pd}{\pd y}\right)
\]
with $h$ holomorphic; if $0$ is an isolated nondegenerate singularity, then $h$ is a unit.
\end{proposition}

\begin{proof}[Proof of Proposition~\ref{prop:cartesian}]


We first record the elementary tangency equation. If
\[
 X=P(x,y)\frac{\pd}{\pd x}+Q(x,y)\frac{\pd}{\pd y},
\]
then $X$ is complex tangent to $S^3(R)$ at $(x,y)$ precisely when
\begin{equation}\label{eq:tangency}
 \ol{x}P(x,y)+\ol{y}Q(x,y)=0.
\end{equation}

Since
\[
 \pd D_1\times\pd D_2\subset S^3(R),
\]
fixing $y_0\in\pd D_2$ shows that $|x|$ is constant on $\pd D_1$, and fixing $x_0\in\pd D_1$ shows that $|y|$ is constant on $\pd D_2$. Hence, for some $a,b>0$ with $a^2+b^2=R^2$,
\[
 D_1=\{|x|<a\},\qquad D_2=\{|y|<b\},
\]
and
\[
 T=\{|x|=a,\ |y|=b\}.
\]
Consider
\[
 Y=b^2x\frac{\pd}{\pd x}-a^2y\frac{\pd}{\pd y}.
\]
Both $X$ and $Y$ are complex tangent to $S^3(R)$ along $T$. Since $T_p^{\C}S^3(R)$ is one-dimensional, $X$ and $Y$ are collinear there. Thus the holomorphic function
\[
 X\wedge Y=-a^2yP-b^2xQ
\]
vanishes on $T$.

A holomorphic function vanishing on the Clifford torus $T$ vanishes identically. Indeed, if
\[
 F(x,y)=\sum_{m,n\ge0}c_{mn}x^my^n,
\]
then on $x=ae^{i\theta}$, $y=be^{i\varphi}$ the vanishing of $F$ and uniqueness of the double Fourier series give $c_{mn}=0$ for all $m,n$. Hence
\[
 a^2yP+b^2xQ\equiv0.
\]
Putting $x=0$ and $y=0$ gives $P=xp$ and $Q=yq$, with $p,q$ holomorphic, and therefore
\[
 a^2p+b^2q=0.
\]
Thus
\[
 P=b^2xh,\qquad Q=-a^2yh
\]
for a holomorphic function $h$. Nondegeneracy gives $h(0)\neq0$, so $h$ is a unit after shrinking the neighborhood. This proves the proposition.
\end{proof}

\section{Proofs of the one-boundary rigidity results}


\begin{proof}[Proof of Theorem~\ref{thm:polynomial}]
Write
\[
 \Phi(z,w)=\sum_{0\le j\le M,\,0\le k\le N}v_{jk}z^jw^k,
 \qquad v_{jk}\in\C^2,
\]
where $(M,N)$ is the bidegree. Since $\Phi(0)=0$, $v_{00}=0$, and since $D\Phi(0)$ is invertible,
\[
 p:=v_{10},\qquad q:=v_{01}
\]
form a basis of $\C^2$. The spherical boundary identity
\begin{equation}\label{eq:sphere}
 |\Phi(z,w)|^2=R^2,\qquad (z,w)\in\T^2,
\end{equation}
gives, for every nonzero $(\alpha,\beta)\in\mathbb Z^2$,
\begin{equation}\label{eq:fourier}
 \sum_{j,k}\ip{v_{j+\alpha,k+\beta}}{v_{jk}}=0,
\end{equation}
with coefficients outside the support understood to vanish.

Define the reflected polynomial vector
\[
 \Phi^{\#}(z,w)=z^Mw^N\,\ol{\Phi(1/\ol z,1/\ol w)}
\]
and the polynomial one-form
\[
 \Omega=\Phi^{\#}\cdot d\Phi.
\]
On $\T^2$ one has
\[
 \Phi^{\#}\cdot\Phi=z^Mw^N|\Phi|^2=R^2z^Mw^N.
\]
Both sides are polynomials, so uniqueness on the totally real torus gives the polynomial identity
\begin{equation}\label{eq:reflected-identity}
 \Phi^{\#}\cdot\Phi=R^2z^Mw^N.
\end{equation}
This identity also controls saturation. If a nonconstant polynomial $d$ divides both components of $\Phi^{\#}$, then~\eqref{eq:reflected-identity} gives $d\mid z^Mw^N$. Thus every irreducible factor of $d$ is either $z$ or $w$. By the definition of the bidegree, at least one component of $\Phi$ contains a monomial of $z$-degree $M$, so at least one component of $\Phi^{\#}$ has a term not divisible by $z$; hence $z$ is not a common factor. Similarly $w$ is not a common factor. Therefore the two components of $\Phi^{\#}$ are relatively prime.

The coefficient vector of $\Omega$ is $(D\Phi)^t\Phi^{\#}$. Since $D\Phi(0)$ is invertible, $D\Phi$ is invertible over the local holomorphic ring at the origin. Consequently the coefficients of $\Omega$ generate the same local ideal as the components of $\Phi^{\#}$, and $\Omega$ is already saturated near $0$.

On $\T^2$, $\Omega$ is a nonzero scalar multiple of the pullback of the complex contact form of the sphere. Since $\Phi^*\F$ defines the same line field there, the wedge of a holomorphic defining form of $\Phi^*\F$ with $\Omega$ vanishes on the maximally totally real torus $\T^2$, hence identically. Thus the foliation defined by $\Omega$ is precisely $\Phi^*\F$ near the origin.

If $v_{MN}\neq0$, then $\Phi^{\#}(0)=\ol{v_{MN}}\neq0$, and therefore $\Omega(0)\neq0$ because $D\Phi(0)$ is invertible. This contradicts the fact that $0$ is a singular point of $\Phi^*\F$. Hence
\[
 v_{MN}=0.
\]
Set
\[
 u=v_{M-1,N},\qquad v=v_{M,N-1}.
\]
These are precisely the coefficients determining the linear part of $\Phi^{\#}$. If $u=v=0$, then $\Phi^{\#}$, and hence the already saturated form $\Omega$, has order at least two at the origin. The induced foliation would therefore have a degenerate singularity, contrary to the hypotheses. Thus at least one of $u,v$ is nonzero, and the linear part of $\Omega$ is
\[
 \Omega_1=(\ol u z+\ol v w)\cdot(p\,dz+q\,dw).
\]
Put
\[
 A=\ip{p}{u},\quad B=\ip{q}{u},\quad
 C=\ip{p}{v},\quad D=\ip{q}{v}.
\]
The Fourier relation~\eqref{eq:fourier} for the mode $(2-M,-N)$ gives $A=0$: indeed, the only potentially nonzero contribution is $\ip{p}{u}$, the two endpoint contributions involving $v_{00}$ and $v_{MN}$ being zero. Similarly the mode $(-M,2-N)$ gives $D=0$. Finally, if $(M,N)\neq(1,1)$, the nonzero mode $(1-M,1-N)$ gives
\[
 B+C=0,
\]
because the only surviving terms are $\ip{q}{u}$ and $\ip{p}{v}$. Consequently
\[
 \Omega_1=B(z\,dw-w\,dz).
\]
Since the singularity is nondegenerate, $B\neq0$. An annihilating vector field has radial linear part
\[
 B\left(z\frac{\pd}{\pd z}+w\frac{\pd}{\pd w}\right),
\]
which is of Poincar\'e type and has eigenvalue ratio $+1$. This contradicts the Siegel hypothesis. Therefore
\[
 (M,N)=(1,1).
\]

In this remaining case $v_{11}=v_{MN}=0$, so
\[
 \Phi(z,w)=pz+qw.
\]
The Fourier mode $(1,-1)$ in~\eqref{eq:fourier} gives $\ip{p}{q}=0$. Since $p$ and $q$ form a basis, after a unitary target change we obtain
\[
 \Phi(z,w)=(az,bw),\qquad a,b\neq0.
\]
Thus $\Phi$ is linear. Moreover the pulled-back characteristic foliation is
\[
 |b|^2z\frac{\pd}{\pd z}-|a|^2w\frac{\pd}{\pd w},
\]
up to multiplication by a holomorphic unit. Hence $\F$ is analytically linearizable.
\end{proof}

\begin{proof}[Proof of Theorem~\ref{thm:logarithmic}]
Choose a smooth defining function $\rho$ for $D$ which is strictly plurisubharmonic in a collar of $\partial D$, and put
\[
 u=\rho\circ\Phi.
\]
Then $u$ is real valued, $u=0$ on $\T^2$, and $u\le0$ on $\ol{\D}^2$ by the hypothesis $\Phi(\ol{\D}^2)\subset\ol D$. Set, on $\T^2$,
\[
 A=z\,u_z,\qquad B=w\,u_w.
\]
Since $u$ is constant on $\T^2$, differentiation in the two angular variables gives
\[
 \Im A=\Im B=0,
\]
so $A$ and $B$ are real there.

We next prove positivity. Fix $w_0\in\T$ and consider the nonconstant holomorphic disc
\[
 \gamma_{w_0}:\D\longrightarrow\C^2,
 \qquad z\longmapsto\Phi(z,w_0).
\]
The function $v(z)=u(z,w_0)$ satisfies $v\le0$ on $\ol\D$ and $v=0$ on $\T$. Near the boundary circle the image of the disc lies in the collar where $\rho$ is strictly plurisubharmonic; hence $v$ is strictly subharmonic there. The boundary Hopf lemma gives a strictly positive outward radial derivative. Thus, for $|z|=1$,
\[
 \frac{d}{dt}\Big|_{t=1}u(tz,w_0)
 =zu_z+\ol z\,u_{\ol z}
 =2\Re A
 =2A>0.
\]
Therefore $A>0$ on $\T^2$. The same argument, fixing $z_0\in\T$ and using the discs $w\mapsto\Phi(z_0,w)$, gives $B>0$.

The $(1,0)$-part of $du$ on $\T^2$ is
\[
 u_z\,dz+u_w\,dw
 =A\frac{dz}{z}+B\frac{dw}{w}.
\]
Hence the pulled-back complex tangent line to $\partial D$ is defined on $\T^2$ by
\[
 A(z,w)\frac{dz}{z}+B(z,w)\frac{dw}{w}=0.
\]
Because $T\subset M(\F,\partial D)$, the logarithmic form
\[
 \omega=a(z,w)\frac{dz}{z}+b(z,w)\frac{dw}{w}
\]
defines the same line there. Consequently
\[
 aB-bA=0,
 \qquad
 \frac{a}{b}=\frac{A}{B}>0
 \quad\text{on }\T^2.
\]

The quotient $H=a/b$ is holomorphic on a neighborhood of $\ol{\D}^2$ and real valued on $\T^2$. For each fixed $w\in\T$, the imaginary part of $z\mapsto H(z,w)$ vanishes on $\T$, hence on $\D$; then fixing $z\in\D$ and repeating the argument in $w$ shows that $H$ is real valued on $\D^2$. A holomorphic real-valued function is constant, so
\[
 H\equiv c>0.
\]
Therefore
\[
 \omega=b(z,w)\left(c\frac{dz}{z}+\frac{dw}{w}\right).
\]
Since $b$ is a unit, $\G$ is the linear Siegel foliation generated by
\[
 z\frac{\pd}{\pd z}-c\,w\frac{\pd}{\pd w}.
\]
Transporting this foliation by $\Phi$ proves that the germ of $\F$ at $0$ is analytically linearizable of Siegel type.
\end{proof}

\section{Weighted-Hopf Models and Transversely Holomorphic Fillings}
In what follows, the word {\em flow} refers to a real flow. We record several facts about nonsingular transversely holomorphic flows on $S^3$, together with a filling construction for linear boundary dynamics. They provide a global framework for the invariant-torus geometry considered later. We begin with the linear Siegel model, which also separates the natural transverse-holomorphic structure inherited from the ambient foliation from alternative transverse-holomorphic structures carried by the desingularized real flow.




\begin{exe}[The linear Siegel model and the three boundary structures]
\label{section:siegellinear}
Consider
\[
 X_0=\lambda z\frac{\pd}{\pd z}-\mu w\frac{\pd}{\pd w},
 \qquad \lambda,\mu>0,
\]
and the round sphere $S^3(r)$. Since
\[
 X_0(|z|^2+|w|^2)=\lambda|z|^2-\mu|w|^2,
\]
the complex tangency set is the Clifford torus
\begin{equation}\label{eq:linear-torus}
 T(r)=\left\{|z|^2=\frac{\mu}{\lambda+\mu}r^2,\quad
 |w|^2=\frac{\lambda}{\lambda+\mu}r^2\right\}.
\end{equation}
On $S^3(r)\setminus T(r)$ the real intersection line $T\F\cap TS^3(r)$ is generated by
\[
 Y=\lambda\frac{\pd}{\pd\theta_z}-\mu\frac{\pd}{\pd\theta_w}.
\]
The same formula defines a nonsingular real-analytic flow on the whole sphere. Moreover, every Hopf torus
\[
 T_c=\{|z|^2=cr^2,\ |w|^2=(1-c)r^2\},\qquad 0<c<1,
\]
is invariant under $Y$; the torus in~\eqref{eq:linear-torus} is only the distinguished member on which the ambient complex foliation is tangent to the sphere. Thus the extended real flow has infinitely many invariant tori.

The natural transverse holomorphic structure induced by the ambient Siegel foliation on $S^3(r)\setminus T(r)$ does not extend through $T(r)$. Indeed, away from the coordinate axes a local logarithmic first integral is
\[
 H=\mu\log z+\lambda\log w.
\]
Writing $s=|z|^2$ on the sphere, the real part of $H$ is
\[
 R(s)=\frac{\mu}{2}\log s+\frac{\lambda}{2}\log(r^2-s),
\]
and
\[
 R'(s)=0
\]
exactly at $s=\mu r^2/(\lambda+\mu)$, i.e. on $T(r)$. Hence the natural transverse complex coordinate loses real rank precisely at the tangency torus.

Nevertheless, the globally extended real flow $Y$ does admit another transverse holomorphic structure. The diffeomorphism
\[
 \kappa(z,w)=(z,\ol w)
\]
conjugates $Y$ to the positive weighted Hopf flow
\[
 \lambda\frac{\pd}{\pd\theta_z}+\mu\frac{\pd}{\pd\theta_w},
\]
so pulling back its standard transverse holomorphic structure gives a global transverse holomorphic structure for $Y$. This example separates three different notions: extension of the underlying real flow, extension of the natural ambient transverse holomorphic structure, and existence of some global transverse holomorphic structure on the extended real flow.

When $\lambda/\mu\in\mathbb Q$, write $\lambda/\mu=p/q$ with coprime positive integers $p,q$. Then the flow of $Y$ is periodic on every invariant Hopf torus $T_c$. Hence Corollary~\ref{cor:local-th-periodic} is realized here by the familiar rational weighted-Hopf, or local Seifert, picture: a neighborhood of the distinguished tangency torus is foliated by closed circles and carries the alternative transversely holomorphic structure coming from the weighted-Hopf model. Thus the periodic tangency-torus result recovers locally, without assuming a first integral, the boundary geometry exhibited by the resonant linear Siegel singularity.

\end{exe}


\begin{lemma}[compact invariant surfaces]\label{lem:compact-surface}
Let $\mathcal L$ be a nonsingular flow on $S^3$. If $\Sigma\subset S^3$ is a compact embedded invariant surface, then $\Sigma$ is diffeomorphic to the 2-torus $T^2$.
\end{lemma}

\begin{proof}
The generating vector field restricts to a nowhere-vanishing tangent vector field on $\Sigma$, so $\chi(\Sigma)=0$ by the Poincar\'e--Hopf theorem. Every closed embedded surface in $S^3$ is orientable. Hence $\Sigma$ is a closed orientable surface of Euler characteristic zero, and therefore $\Sigma\simeq T^2$.
\end{proof}

For comparison with the boundary-rigidity method of Brunella in \cite{Brunella1995}, we also record the following reconstruction mechanism in the everywhere-transverse situation.

\begin{lemma}[weighted-Hopf boundary reconstruction]\label{lem:hopf-boundary}
Let $\mathcal H$ be a nonsingular holomorphic foliation on a neighborhood of the sphere $S^3(R)$ which is everywhere transverse to $S^3(R)$. Assume that the induced transversely holomorphic flow on $S^3(R)$ is transversely holomorphically conjugate to the weighted Hopf flow induced by
\[
 Z=\lambda z\frac{\pd}{\pd z}+\mu w\frac{\pd}{\pd w},
 \qquad \lambda,\mu>0.
\]
Then, after shrinking the neighborhood of $S^3(R)$, $\mathcal H$ is defined by a closed meromorphic one-form. Its polar divisor is contained in the complex leaves through the two closed Hopf orbits corresponding, in the model, to $\{z=0\}\cap S^3(R)$ and $\{w=0\}\cap S^3(R)$.
\end{lemma}

\begin{proof}
The linear model is defined by the closed meromorphic form
\[
 \omega_0=\mu\frac{dz}{z}-\lambda\frac{dw}{w},
 \qquad d\omega_0=0.
\]
Away from the two coordinate Hopf circles, the restriction of $\omega_0$ to the sphere is a basic transversely holomorphic one-form for the induced flow. Pulling it back by a transverse-holomorphic conjugacy gives a basic transversely holomorphic one-form $\alpha$ on the complement of the corresponding two closed orbits.

Let $(x,\zeta)$ be a holomorphic flow box for $\mathcal H$, so that the leaves are given by $\zeta=\mathrm{constant}$. Since $\mathcal H$ is transverse to the sphere, $\zeta|_{S^3(R)}$ is a genuine local transverse complex coordinate, and therefore
\[
 \alpha=a(\zeta)\,d\zeta
\]
with $a$ holomorphic. Thus $\alpha$ extends canonically to the ambient flow box as the closed holomorphic form $a(\zeta)d\zeta$. On overlaps, two such extensions agree on an open subset of the real-analytic hypersurface $S^3(R)$, hence agree identically by holomorphic uniqueness. They therefore glue to a closed holomorphic defining form near the sphere away from the two closed orbits.

Near either closed orbit, the transverse-holomorphic conjugacy identifies its holonomy germ with that of the corresponding Hopf orbit. We use here the standard suspension-extension argument in the form employed by Brunella: equality of the holomorphic holonomy germs yields a biholomorphic conjugacy between sufficiently small ambient neighborhoods of the corresponding closed boundary orbits; see~\cite[pp.~124--125]{Brunella1995}. Pulling back $\omega_0$ by this local biholomorphism gives a closed meromorphic defining form with a simple polar component along the complex leaf through the orbit. Choosing the suspension to agree with the given transverse conjugacy fixes the transverse normalization, so these local meromorphic forms glue with the form already constructed off the two closed orbits. This gives the desired closed meromorphic one-form on a full neighborhood of $S^3(R)$.
\end{proof}

\begin{remark}
Brunella's Hartogs--Levi extension step applies once a closed meromorphic defining form has been obtained on a neighborhood of the whole boundary; see~\cite[pp.~124--125]{Brunella1995}. Lemma~\ref{lem:hopf-boundary} concerns the everywhere-transverse setting. In the tangency-torus situations of Theorems~\ref{thm:main} and~\ref{thm:spherical-characterization}, the natural transverse holomorphic structure degenerates along the tangency torus, so the lemma cannot be applied directly.
\end{remark}

\begin{lemma}[invariant torus and weighted Hopf dynamics]\label{lem:torus-hopf}
Let $\mathcal L$ be a nonsingular transversely holomorphic flow on $S^3$. If $\mathcal L$ admits an invariant embedded two-torus, then $\mathcal L$ is transversely holomorphically conjugate to a weighted Hopf flow.
\end{lemma}

\begin{proof}
This is a direct consequence of the Brunella--Ghys classification of transversely holomorphic flows, together with the explicit $S^3$ models; see~\cite{BrunellaGhys1995,Ghys1996} and, for an explicit exhaustive list on $S^3$,~\cite[Theorem~4.9 and Section~6]{GeigesGonzalo2016}. In the explicit classification, the nonreal members of the continuous Poincar\'e family have only the two Hopf circles as compact leaves and every other leaf is asymptotic to them, while the discrete family has only one compact leaf and the remaining leaves have the corresponding asymptotic behavior. Neither case admits a compact invariant torus. The remaining real family is precisely the weighted Hopf family, whose leaves lie on the Hopf tori $\{|z_1|=\mathrm{const.}\}$.
\end{proof}

\begin{corollary}\label{cor:many-tori}
If a nonsingular transversely holomorphic flow on $S^3$ admits one compact embedded invariant surface, then it is transversely holomorphically conjugate to a weighted Hopf flow. Consequently it admits a one-parameter family of invariant tori filling the complement of two distinguished closed orbits.
\end{corollary}

\begin{proof}
By Lemma~\ref{lem:compact-surface}, the invariant surface is a torus, and Lemma~\ref{lem:torus-hopf} applies. In the weighted Hopf model the invariant tori are
\[
 T_c=\{|z_1|^2=c,\ |z_2|^2=1-c\},\qquad 0<c<1,
\]
and they fill the complement of the two Hopf circles.
\end{proof}

\begin{corollary}[compact transversely holomorphic flows on $S^3$]\label{cor:compact-th}
If every leaf of a nonsingular transversely holomorphic flow on $S^3$ is compact, then the flow is transversely holomorphically conjugate to a rational weighted Hopf flow. Equivalently, after a constant rescaling of the generator, it has the model
\[
 p\frac{\pd}{\pd\theta_1}+q\frac{\pd}{\pd\theta_2},
 \qquad p,q\in\mathbb N,\quad \gcd(p,q)=1.
\]
In particular, it defines a Seifert fibration of $S^3$.
\end{corollary}

\begin{proof}
By the classification quoted in Lemma~\ref{lem:torus-hopf}, the compact-leaf case belongs to the real weighted Hopf family. The explicit description in~\cite[Propositions~6.1 and~6.2]{GeigesGonzalo2016} shows that all leaves are closed exactly when the weight ratio is rational; the resulting foliation is a Seifert fibration.
\end{proof}

\begin{remark}
Corollaries~\ref{cor:many-tori} and~\ref{cor:compact-th} are recorded only for convenience; their content is contained in the established classification theory and the explicit $S^3$ models.
\end{remark}

\begin{lemma}[linear characteristic-flow filling]\label{lem:linear-filling}
Let $V\simeq S^1\times\D$ be a solid torus with boundary $T=\partial V$, and fix a real-analytic product identification whose boundary coordinates $(\theta,\varphi)$ are longitude--meridian coordinates. Suppose that, in these boundary coordinates, a nonsingular real-analytic foliation on $T$ is real-analytically conjugate by a boundary diffeomorphism extending to a real-analytic diffeomorphism of $V$ to the constant linear foliation generated by
\[
 a\frac{\pd}{\pd\theta}+b\frac{\pd}{\pd\varphi},\qquad a\ne0,
\]
where $\varphi$ is meridional. Then the boundary foliation admits a nonsingular transversely holomorphic extension to $V$. More precisely, it is the boundary of the suspension
\[
 V_\alpha=(\D\times[0,1])/((z,1)\sim(e^{i\alpha}z,0)),
 \qquad \alpha=2\pi\frac ba,
\]
and the core circle is a closed leaf with rotational holonomy $z\mapsto e^{i\alpha}z$.
\end{lemma}

\begin{proof}
The circle $\{\theta=0\}$ is a global transversal for the linear boundary flow. After one turn in the longitudinal coordinate, the meridional coordinate changes by $2\pi b/a$, so the first-return map is the rigid rotation $e^{i\varphi}\mapsto e^{i(\varphi+2\pi b/a)}$. Suspend the holomorphic disk automorphism $z\mapsto e^{i\alpha}z$. The vertical foliation of $V_\alpha$ has holomorphic transverse coordinate $z$, and its only transverse gluing map is the holomorphic rotation $z\mapsto e^{i\alpha}z$. Hence it is transversely holomorphic. The fixed point $z=0$ suspends to the core closed leaf. Pulling the construction back by the assumed real-analytic extension of the boundary conjugacy to $V$ gives the asserted filling.
\end{proof}

\begin{corollary}[periodic filling]\label{cor:periodic-filling}
Under the hypotheses of Lemma~\ref{lem:linear-filling}, if $b/a\in\mathbb Q$, then the core holonomy is finite. In particular, if an invariant torus $T\subset S^3$ splits $S^3=V_1\cup_TV_2$ into two solid tori and the induced linear boundary foliations admit compatible fillings on both sides, then the resulting global transversely holomorphic flow on $S^3$ is of rational weighted-Hopf type, hence a Seifert fibration.
\end{corollary}

\begin{proof}
If $b/a=q/p$ in lowest terms, then the disk rotation has order $p$. Compatible fillings give a nonsingular transversely holomorphic flow on $S^3$ with invariant torus $T$. Lemma~\ref{lem:torus-hopf} gives weighted-Hopf dynamics, and the rationality of the rotation number gives the rational weighted-Hopf, equivalently Seifert, case.
\end{proof}


\section{General Geometry of Smooth Tangency Tori}
The preceding rigidity results concern situations in which the geometry of the tangency set forces linearization. We conclude by considering a smooth tangency torus without imposing those global rigidity hypotheses. The rank-one identity of Lemma~\ref{lem:spherical-rank}, used earlier to prove automatic spectral alignment in the spherical theorem, is also the starting point here. The principal result of this section shows that periodicity of the characteristic foliation forces a stronger local rigidity: the characteristic dynamics can be desingularized across the torus and both the holomorphic and real return holonomies are trivial.

\begin{theorem}[periodic tangency-torus rigidity]\label{thm:periodic-torus-rigidity}
Let $\F$ be a one-dimensional holomorphic foliation regular in a neighborhood of $S^3(r)\subset\C^2$, and let $T$ be a smooth embedded two-torus which is a connected component of the complex tangency locus $M(\F,S^3(r))$. Assume that the characteristic foliation induced on $T$ is a circle fibration. Then:
\begin{enumerate}
\item on a neighborhood of $T$ in $S^3(r)$, the characteristic line field on $S^3(r)\setminus T$ extends to a nonsingular real-analytic vector field $Y$ tangent to the real two-planes underlying the complex leaves of $\F$; the field $Y$ leaves $T$ invariant and restricts on $T$ to its characteristic circle foliation;
\item for every characteristic circle $\gamma\subset T$, the holomorphic holonomy germ of $\gamma$ in $\F$ is the identity;
\item the Poincar\'e return germ of $Y$ along every such $\gamma$ is the identity.
\end{enumerate}
\end{theorem}

\begin{proof}
\emph{Step 1: desingularization.}
Let
\[
 X=P_1(z)\frac{\pd}{\pd z_1}+P_2(z)\frac{\pd}{\pd z_2}
\]
be a holomorphic vector field generating $\F$ near $T$, and put
\[
 \rho=|z_1|^2+|z_2|^2-r^2,\qquad
 g=X(\rho)|_{S^3(r)}=P_1\ol z_1+P_2\ol z_2.
\]
By Lemma~\ref{lem:spherical-rank},
\begin{equation}\label{eq:barXg}
 \ol X(g)=|P_1|^2+|P_2|^2=\|X\|^2>0
\end{equation}
along $T$, and the real differential of $g:S^3(r)\to\C$ has rank one there.

Although $T$ was assumed only smooth, this rank-one property forces it to be real analytic. Indeed, near each $p\in T$ one real component of the real-analytic map $g:S^3(r)\to\C$ has nonzero differential, so its zero set is a real-analytic surface containing $T$. Since both have real dimension two, the inclusion of $T$ into that surface is locally open; hence $T$ agrees locally with a real-analytic hypersurface. Thus $T$ is a compact real-analytic two-sided hypersurface of $S^3(r)$. By the real-analytic tubular-neighborhood theorem, after shrinking around $T$ there is a global real-analytic defining function $\sigma$ on a neighborhood $U$ such that
\[
 T=\{\sigma=0\},\qquad d\sigma\ne0.
\]
Since $g$ vanishes precisely to first order in the normal direction, real-analytic division gives
\begin{equation}\label{eq:g-sigma-h}
 g=\sigma h
\end{equation}
with $h:U\to\C$ real analytic and nowhere zero. Define
\begin{equation}\label{eq:Ydesing}
 Y=\Re(i\ol h\,X)|_{S^3(r)}.
\end{equation}
Then $Y$ is real analytic and nonsingular on $U$. Moreover,
\[
 Y(\rho)=\Re(i\ol h\,X(\rho))
          =\Re(i\sigma|h|^2)=0,
\]
so $Y$ is tangent to $S^3(r)$, and by construction it is tangent to the real plane underlying the complex leaf of $\F$. Off $T$, the standard characteristic generator
\[
 V=\Re(i\ol g\,X)
\]
satisfies $V=\sigma Y$, so $Y$ extends the same characteristic line field.

It remains to prove that $T$ is $Y$-invariant. On $T$, equation~\eqref{eq:g-sigma-h} gives
\[
 dg(\xi)=h\,d\sigma(\xi),\qquad \xi\in TS^3(r),
\]
and hence $dg(\xi)/h\in\R$. Since at a tangency point both $\Re X$ and $\Im X$ are tangent to the sphere, put
\[
 A=Xg,\qquad S=\ol Xg=\|X\|^2>0.
\]
Applying the preceding reality condition to $\Re X$ and $\Im X$ yields
\[
 \frac{A+S}{h}\in\R,
 \qquad
 \frac{A-S}{ih}\in\R.
\]
If $a=A/h$ and $b=S/h$, these say that $a+b$ is real and $a-b$ is purely imaginary; hence $a=\ol b$. Since $S$ is real,
\begin{equation}\label{eq:Xg-relation}
 A=\frac{h}{\ol h}\,S.
\end{equation}
Therefore, on $T$,
\[
 Y(g)=\frac{i}{2}(\ol h\,Xg-h\,\ol Xg)
     =\frac{i}{2}\left(\ol h\,\frac{h}{\ol h}S-hS\right)=0.
\]
Using $g=\sigma h$ and $\sigma=0$ on $T$, we obtain
\[
 0=Y(g)=h\,Y(\sigma),
\]
so $Y(\sigma)=0$ on $T$. Thus $Y$ is tangent to $T$, and its restriction there is precisely the characteristic line field.

\smallskip
\noindent\emph{Step 2: trivial holonomy.}
Fix $p\in\gamma$ and a small complex transversal $\Delta$ to the complex leaf of $\F$ through $\gamma$. Since $T$ is foliated by characteristic circles, $I=T\cap\Delta$ is, after shrinking, a real-analytic transverse arc and nearby points of $I$ lie on nearby characteristic circles of the same circle fibration. The holonomy germ $h_\gamma:(\Delta,p)\to(\Delta,p)$ associated with one turn around $\gamma$ therefore fixes every point of $I$. By the one-variable holomorphic identity theorem,
\[
 h_\gamma\equiv\operatorname{id}.
\]
Trivial holonomy around $\gamma$ gives, after shrinking to a foliated tubular neighborhood of $\gamma$, a single-valued holomorphic transverse coordinate
\[
 F:W\longrightarrow\Delta
\]
constant on the plaques of $\F$.

Let $\Sigma\subset S^3(r)$ be a small real-analytic Poincar\'e section to the nonsingular field $Y$ through $p$, and let $P:(\Sigma,p)\to(\Sigma,p)$ be the return germ corresponding to one turn around $\gamma$. Because $Y$ is tangent to the complex leaves,
\[
 F\circ P=F.
\]
On either side of $T$, the holomorphic foliation is transverse to $S^3(r)$; hence $F|_\Sigma$ has real rank two there and is locally a real-analytic diffeomorphism. It follows that $P=\operatorname{id}$ on a nonempty open subset of one side of $\Sigma\setminus(\Sigma\cap T)$. Since $P$ is real analytic, the real-analytic identity principle gives
\[
 P\equiv\operatorname{id}
\]
as a germ at $p$. In particular $DP_p=I$.
\end{proof}

\begin{corollary}[local TH structure in the periodic case]\label{cor:local-th-periodic}
Under the hypotheses of Theorem~\ref{thm:periodic-torus-rigidity}, there exists, after shrinking around $T$, a saturated neighborhood
\[
 U\subset S^3(r)
\]
of $T$ on which the desingularized field $Y$ defines a regular real-analytic circle fibration. In particular, the orbit foliation of $Y|_U$ admits a transversely holomorphic structure extending across $T$.
\end{corollary}

\begin{proof}
By Theorem~\ref{thm:periodic-torus-rigidity}, every characteristic circle $\gamma\subset T$ has trivial Poincar\'e return germ for $Y$. The local stability theorem for a compact leaf with finite holonomy therefore gives, for each such $\gamma$, a saturated neighborhood $U_\gamma$ in which every leaf of $Y$ is a circle and the foliation is locally a product
\[
 S^1\times D^2\longrightarrow D^2.
\]
Since $T$ is compact, finitely many such neighborhoods cover it. After shrinking their union, we obtain a saturated neighborhood $U$ of $T$ all of whose leaves are compact circles. The holonomy is trivial near $T$, so the local leaf spaces glue without orbifold isotropy and define a real-analytic surface $B$ for which the quotient map
\[
 \pi:U\longrightarrow B
\]
is a real-analytic circle fibration. Since $T$ is two-sided in $S^3(r)$, $U$ may be chosen inside a tubular neighborhood diffeomorphic to $T^2\times(-\varepsilon,\varepsilon)$.

The oriented normal bundle of the oriented flow $Y$ is of real rank two, hence the local leaf space $B$ is oriented. Choose a real-analytic conformal structure on $B$; equivalently, regard $B$ locally as a Riemann surface. Pulling back its complex coordinate charts by $\pi$ gives transverse complex coordinates for the orbit foliation of $Y|_U$, with holomorphic transition functions. Thus $Y|_U$ is transversely holomorphic across the tangency torus.
\end{proof}

\begin{remark}\label{rem:periodic-local-model}
Corollary~\ref{cor:local-th-periodic} gives a local resolution of the transverse-holomorphic degeneracy at a periodic tangency torus. The transverse holomorphic structure naturally inherited from the ambient holomorphic foliation need not extend through $T$, but the desingularized real characteristic flow nevertheless carries an alternative transverse holomorphic structure on a full collar of $T$. No global compatibility with the natural transverse structures on the two complementary regions of the sphere is asserted here.
\end{remark}

\begin{thebibliography}{9}

\bibitem{Arnold}
V.~I. Arnold,
\emph{Chapitres suppl\'ementaires de la th\'eorie des \'equations diff\'erentielles ordinaires},
Mir, 1980.

\bibitem{Brjuno}
A.~D. Brjuno,
Analytical form of differential equations,
\emph{Trans. Moscow Math. Soc.} \textbf{25} (1971), 131--288.

\bibitem{Brunella1995}
M.~Brunella,
\emph{On holomorphic foliations in complex surfaces transverse to a sphere},
Bol. Soc. Brasil. Mat. (N.S.) \textbf{26} (1995), no.~2, 117--128.

\bibitem{BrunellaI}
Brunella, M.: On transversely holomorphic flows I.
Invent. Math. \textbf{126}, 265--279 (1996).

\bibitem{BrunellaGhys1995}
M.~Brunella and E.~Ghys,
\emph{Umbilical foliations and transversely holomorphic flows},
J. Differential Geom. \textbf{41} (1995), no.~1, 1--19.

\bibitem{Camacho-LN} César Camacho and Alcides Lins Neto, Geometric Theory of Foliations, Birkhäuser Boston, Boston, MA, 1985. First edition, VIII + 206 pages.

\bibitem{GeigesGonzalo2016}
H.~Geiges and J.~Gonzalo P\'erez,
\emph{Transversely holomorphic flows and contact circles on spherical $3$-manifolds},
Enseign. Math. \textbf{62} (2016), no.~3--4, 527--567.

\bibitem{Ghys1996}
E.~Ghys,
\emph{On transversely holomorphic flows. II},
Invent. Math. \textbf{126} (1996), no.~2, 281--286.

\bibitem{Haefliger} André Haefliger, Variétés feuilletées,
Annali della Scuola Normale Superiore di Pisa, Série III, 16 (1962), no. 4, 367–397.

\bibitem{Ilyashenko-Yakovenko} Yulij Ilyashenko and Sergei Yakovenko, Lectures on Analytic Differential Equations, Graduate Studies in Mathematics, Vol. 86, American Mathematical Society, Providence, Rhode Island, 2008.

\bibitem{Ito} Toshikazu Ito, A Poincaré–Bendixson type theorem for holomorphic vector fields, Sūrikaisekikenkyūsho Kōkyūroku, No. 878 (1994), pp. 1–9. First  appeared in the proceedings Singularities of holomorphic vector fields and related topics (Kyoto, 1993).

\bibitem{Ito1996} Toshikazu Ito, The number of compact leaves of a one-dimensional foliation on the $2n-1$ dimensional sphere $S^{2n-1}$ associated with a holomorphic vector field, Sūrikaisekikenkyūsho Kōkyūroku, No. 955 (1996), pp. 75--79.

\bibitem{MartinetRamis}
J. Martinet and J.-P. Ramis,
Classification analytique des \'equations diff\'erentielles non lin\'eaires r\'esonantes du premier ordre,
\emph{Ann. Sci. \'Ecole Norm. Sup. (4)} \textbf{16} (1983), 571--621.

\bibitem{MatteiMoussu}
J.-F. Mattei and R. Moussu,
Holonomie et int\'egrales premi\`eres,
\emph{Ann. Sci. \'Ecole Norm. Sup. (4)} \textbf{13} (1980), 469--523.

\end{thebibliography}

\bigskip
\noindent
Toshikazu Ito\\
Professor Emeritus, Ryukoku University\\
67 Tsukamoto-cho, Fukakusa, Fushimi-ku, Kyoto 612-8577, Japan

\medskip
\noindent
Bruno Sc\'ardua\\
Instituto de Matem\'atica, Universidade Federal do Rio de Janeiro\\
Av. Athos da Silveira Ramos, 149, Centro de Tecnologia, Bloco C\\
Cidade Universit\'aria, Rio de Janeiro, RJ 21941-909, Brazil\\
E-mail: \texttt{bruno.scardua@gmail.com}

\end{document}
